\documentclass[11pt]{article}

% =====================================================================
% A design/tutorial write-up of the Parthenon loop abstraction
% (src/loop_abstraction). Draft 1.
%
% The companion TikZ figure lives beside this file (halo_tikz.tex) and is
% pulled in with \input in the Halos section.
% =====================================================================

\usepackage[margin=1in]{geometry}
\usepackage{amsmath, amssymb}
\usepackage{tikz}
\usetikzlibrary{arrows.meta, positioning, fit, backgrounds}
\usepackage{listings}
\usepackage{xcolor}
\usepackage{booktabs}
\usepackage{hyperref}
\hypersetup{colorlinks=true, linkcolor=blue!55!black, urlcolor=blue!55!black}

% ---- code listing style ----
\definecolor{codebg}{gray}{0.97}
\definecolor{codekw}{rgb}{0.30,0.20,0.60}
\definecolor{codecm}{rgb}{0.30,0.45,0.30}
\definecolor{codestr}{rgb}{0.60,0.20,0.10}
\lstdefinestyle{cpp}{
  language=C++,
  backgroundcolor=\color{codebg},
  basicstyle=\ttfamily\footnotesize,
  keywordstyle=\color{codekw}\bfseries,
  commentstyle=\color{codecm}\itshape,
  stringstyle=\color{codestr},
  showstringspaces=false,
  breaklines=true,
  columns=fullflexible,
  keepspaces=true,
  frame=single,
  framesep=4pt,
  rulecolor=\color{gray!40},
  morekeywords={constexpr,KOKKOS_LAMBDA,auto,Real,template,typename,requires,concept}
}
\lstset{style=cpp}

\newcommand{\code}[1]{\lstinline[basicstyle=\ttfamily]|#1|}
\newcommand{\S}{\mathcal{S}}

\title{\bf The Parthenon Loop Abstraction:\\
A Portable, Compile-Time-Selectable Loop Structure\\
for Block-Structured Kernels}
\author{Parthenon developers}
\date{Draft --- \today}

\begin{document}
\maketitle

\begin{abstract}
Block-structured finite-volume kernels have a portability problem that is not
solved by a portable parallel-for alone: the \emph{loop structure} that performs
well on a CPU is genuinely different from the one that performs well on a GPU,
even when the arithmetic is identical. The Parthenon loop abstraction
(\code{src/loop_abstraction}) lets a kernel be written once against a logical
$(\text{block}, k, j, i)$ iteration space and then be compiled into one of
several concrete loop structures, selected by compile-time tags. This document
motivates the abstraction with a concrete reconstruct-then-Riemann flux kernel,
shows how a performance engineer would hand-write that kernel differently for CPU
and GPU, and then develops the abstraction's vocabulary --- index spaces, loop
tags, inner tags, backends, scratch, halos, and pack views --- as the machinery
that expresses both hand-written forms from a single source.
\end{abstract}

\tableofcontents

% =====================================================================
\section{Introduction}
% =====================================================================

Parthenon kernels iterate over a logical index space of cells indexed by
$(\text{block}, k, j, i)$. The existing \code{par_for} / \code{par_for_outer}
wrappers make the \emph{dispatch} portable: the same call compiles to a flat
Kokkos range on a GPU and to a nested loop on a CPU. What they do not abstract is
the \emph{shape} of the computation --- how work is chunked, where intermediate
results are stored versus recomputed, and where the boundaries between levels of
parallelism fall. For simple pointwise kernels this does not matter. For the
reconstruction/flux kernels at the heart of a Godunov hydro scheme it matters a
great deal, and the best shape is not the same on both architectures.

The loop abstraction exists to close that gap. A kernel declares \emph{what}
index space it covers and \emph{how} its variables are laid out; a small set of
compile-time tags then selects the concrete loop structure. The intent is that a
single kernel body can be compiled into either a CPU-friendly hierarchical form
that stores and reuses intermediate results, or a GPU-friendly flat form that
recomputes them, with no change to the physics code.

\paragraph{Why this is not a reinvention of Kokkos.} Kokkos is a general
performance-portability layer: it can express essentially any parallel pattern,
and it makes very few assumptions about what a kernel does. That generality is a
strength, but it also means the programmer is responsible for choosing and hand-
coding the loop structure, and a generic pattern gives the compiler and library
little to exploit. The loop abstraction takes the opposite stance. It is
deliberately \emph{narrower} --- it targets block-structured
$(\text{block}, k, j, i)$ finite-volume kernels and a handful of loop shapes that
are known to perform well for them --- and it is built \emph{on top of} Kokkos,
not in place of it. The abstraction still dispatches through Kokkos (or a raw host
nest) underneath; what it adds is a vocabulary specialized enough that the library
can make the structural decisions a programmer would otherwise make by hand:
where to chunk, what to store versus recompute, which index is contiguous, how
scratch is sized. The trade-off is explicit: because the abstraction only knows
how to express a constrained family of kernels, it cannot describe an arbitrary
computation the way raw Kokkos can. In exchange, within that family it is easier
to write a kernel that is already close to optimal on both architectures, and to
switch between structures without rewriting the kernel. It is a domain-specific
convenience layer over Kokkos, not a competitor to it.

This document is organized around that claim. Section~\ref{sec:motivation}
develops a single motivating kernel and writes it out by hand twice, once for each
architecture, to make the differences concrete. Section~\ref{sec:onekernel} shows
the same kernel written once in the abstraction and recovers both hand-written
forms by changing a tag. The remaining sections
(\ref{sec:indexspace}--\ref{sec:views}) develop each piece of the vocabulary in
turn, and Section~\ref{sec:halos} treats halos --- the mechanism for
producer/consumer reuse --- in detail, including their behavior in reduced
dimensionality.

% =====================================================================
\section{A motivating example: a finite-volume flux kernel}
\label{sec:motivation}
% =====================================================================

Consider the flux computation in a directionally-split Godunov scheme, for the
$x_1$ sweep. For each cell we need the flux $F_{i-1/2}$ on its lower $x_1$ face.
The standard recipe is:
\begin{enumerate}
\item \textbf{Reconstruct} the left and right interface states at each face
  $i-1/2$ from a stencil of cell-centered values --- e.g.\ for a three-point
  reconstruction, $q_{i-1}, q_i, q_{i+1}$ contribute to the states at face
  $i-1/2$. Call the resulting one-sided states $q^{+}_{i-1/2}$ (extrapolated from
  the left) and $q^{-}_{i-1/2}$ (extrapolated from the right).
\item \textbf{Solve a Riemann problem} at each face from the two one-sided states:
  $F_{i-1/2} = \mathcal{R}\!\left(q^{+}_{i-1/2}, q^{-}_{i-1/2}\right)$.
\end{enumerate}

The key structural observation is that \emph{a reconstructed state is used
twice}. The state $q^{\pm}_{i-1/2}$ produced at face $i-1/2$ feeds the Riemann
solve at that face; but the reconstruction machinery that produces the states at
face $i-1/2$ also produces (or shares most of its work with) the states at the
neighboring face $i+1/2$. Whether it is worth \emph{storing} a reconstructed state
so the neighbor can reuse it, or simply \emph{recomputing} it, is precisely the
decision that differs between CPU and GPU. The next two subsections write the
kernel out both ways.

Throughout, $q$ is a cell-centered field, and we compute the $x_1$ flux
$F$ on $x_1$-faces.

% ---------------------------------------------------------------
\subsection{How you would write it for a CPU}
\label{sec:cpu}
% ---------------------------------------------------------------

On a CPU the dominant concerns are (i) keeping the innermost loop long,
unit-stride, and vectorizable, and (ii) not repeating expensive arithmetic. Both
push toward a \emph{hierarchical} structure: an outer loop that walks planes or
pencils of the block, and an inner loop that sweeps contiguously along $x_1$.
Reconstruction is done \emph{once per face} and stored in a scratch buffer; the
flux loop then reads two neighboring scratch entries per face.

\begin{lstlisting}
// CPU-style: reconstruct once into scratch, reuse for the flux.
for (int k = kb.s; k <= kb.e; ++k) {
  for (int j = jb.s; j <= jb.e; ++j) {

    // scratch spans [is-1, ie]: one extra face on the low side, because the
    // flux at face i reads the reconstructed state at face i-1 as well.
    Real *qp = scratch_plus_row(k, j);   // q^+ at each face in this row
    Real *qm = scratch_minus_row(k, j);  // q^-

    // Producer: reconstruct every face in the (extended) row exactly once.
    const Real *q_row = &q(k, j, 0);
    #pragma omp simd
    for (int i = is - 1; i <= ie; ++i) {
      qp[i] = reconstruct_plus (q_row[i-1], q_row[i], q_row[i+1]);
      qm[i] = reconstruct_minus(q_row[i-1], q_row[i], q_row[i+1]);
    }

    // Consumer: one Riemann solve per face, reading neighboring scratch.
    Real *F_row = &F(k, j, 0);
    #pragma omp simd
    for (int i = is; i <= ie; ++i) {
      F_row[i] = riemann(qp[i-1], qm[i]);   // <-- reuse of qp at face i-1
    }
  }
}
\end{lstlisting}

Three features matter and will reappear as abstraction concepts:
\begin{itemize}
\item \textbf{Two loops over nearly the same range.} The producer runs over
  $[i_s-1, i_e]$ and the consumer over $[i_s, i_e]$. The producer's range is the
  consumer's range \emph{plus one extra face on the low side}, because
  \code{riemann(qp[i-1], qm[i])} reads a produced value at $i-1$. That ``one
  extra face'' is exactly what a \emph{halo} will name (Section~\ref{sec:halos}).
\item \textbf{Scratch reuse.} Each reconstructed state is computed once and read
  by (up to) two Riemann solves. The scratch buffer is sized to the extended
  producer range.
\item \textbf{Contiguous inner loops.} Both inner loops are unit-stride in $i$
  and marked for SIMD; \code{q_row}, \code{qp}, \code{qm} are raw pointers so the
  compiler sees \code{a[i]}-style access. The row shown here is one $i$-pencil, but
  nothing forbids fusing several adjacent $j$- (or $k$-) rows into one longer
  contiguous run to give the vectorizer more to chew on. Because the allocated
  block is contiguous in memory, such a fused run can be walked as a single span;
  the cost is a little extra ghost work where the run crosses row boundaries. The
  abstraction exposes this as a tunable chunk size (Section~\ref{sec:looptags});
  the point to plant now is that ``how much to fuse'' is a knob, not a fixed
  choice.
\end{itemize}

% ---------------------------------------------------------------
\subsection{How you would write it for a GPU}
\label{sec:gpu}
% ---------------------------------------------------------------

On a GPU the concerns invert. Floating-point work is cheap relative to memory
traffic and synchronization, and a single flat grid of threads gives the
scheduler the most freedom to hide latency. Storing reconstructed states in
shared memory would force a team structure and a barrier between the producer and
consumer phases; it is usually cheaper to give each face its own thread and let it
\emph{recompute} whatever reconstructed states it needs, reading the input field
directly through the pack.

\begin{lstlisting}
// GPU-style: one thread per face, recompute reconstruction, no scratch.
par_for(is, ie, jb.s, jb.e, kb.s, kb.e,
  KOKKOS_LAMBDA(const int k, const int j, const int i) {
    // This face needs q^+ from face i-1 and q^- from face i.
    const Real qp_im1 = reconstruct_plus (q(k,j,i-2), q(k,j,i-1), q(k,j,i));
    const Real qm_i   = reconstruct_minus(q(k,j,i-1), q(k,j,i), q(k,j,i+1));
    F(k, j, i) = riemann(qp_im1, qm_i);
  });
\end{lstlisting}

Here \code{reconstruct_plus} at face $i-1$ is computed independently by the thread
that owns face $i-1$ as well --- the value is produced twice, once per consuming
face. That redundant arithmetic is deliberate: it buys a flat, barrier-free,
single-level parallel loop with direct field access and no shared-memory
bookkeeping.

% ---------------------------------------------------------------
\subsection{What differs, and what does not}
\label{sec:differs}
% ---------------------------------------------------------------

Comparing the two, the \emph{physics} --- \code{reconstruct_plus},
\code{reconstruct_minus}, \code{riemann}, and the fact that the flux at a face
reads $q^+$ from the face below and $q^-$ from the face itself --- is
byte-for-byte identical. What differs is entirely structural:

\begin{center}
\begin{tabular}{@{}lll@{}}
\toprule
 & \textbf{CPU form (\S\ref{sec:cpu})} & \textbf{GPU form (\S\ref{sec:gpu})} \\
 & \code{bvoi} / \code{bovi} & \code{boiv} \\
\midrule
Parallelism   & hierarchical (team over planes, & flat (one thread per face) \\
              & SIMD along $i$)                 &  \\
Intermediate  & stored in scratch, reused       & recomputed per consuming face \\
Access        & raw pointers, \code{a[i]}       & direct pack access \\
Barrier       & producer/consumer split         & none \\
\bottomrule
\end{tabular}
\end{center}

An abstraction that lets us write the shared physics once and choose the column at
compile time is exactly what Section~\ref{sec:onekernel} delivers. At the level of
this section the choice is just the one made above: a \emph{hierarchical} form
that chunks the block and stores intermediates in scratch, versus a \emph{flat
pointwise} form that gives each cell its own thread and recomputes. Both are
selected by a single template argument on the \code{IndexSpace}, and the ``one
extra face'' the producer covers becomes a \code{minus_i} \emph{halo}. The named
tags that spell these choices, and the finer distinctions among the hierarchical
variants, are introduced in Section~\ref{sec:onekernel} alongside the code and
developed in Section~\ref{sec:looptags}.

% =====================================================================
\section{The same kernel, written once}
\label{sec:onekernel}
% =====================================================================

Here is the flux kernel expressed in the abstraction. It is written against an
\code{IndexSpace} whose tags are left as template parameters; instantiating those
tags selects the CPU or GPU structure from Section~\ref{sec:motivation} without
touching the body.

\begin{lstlisting}
namespace la = parthenon::loop_abstraction;

// recon_halo names the extra produced faces the flux loop will read: the flux at
// face i reads the reconstructed q^+ at face i-1, i.e. one cell in -i.
using recon_halo = la::halo::minus_i_t;

// F1-face index space over the interior, backend chosen by the tags.
la::IndexSpace<LOOP_TAG, INNER_TAG> idx_space(
    ninner, IndexDomain::interior, /*halo=*/0, pack.GetNBlocks(),
    md, TopologicalElement::F1);
idx_space.AddPerPointScratch<Real, recon_halo>(2);   // q^+ and q^- per face

const auto dx1 = idx_space.GetDelta(X1DIR);           // logical +1 step in x1

la::outer(idx_space, KOKKOS_LAMBDA(const la::InnerIndexRange<decltype(idx_space)> &r,
                                   int b) {
  // Extend this slice by recon_halo: the producer below now covers the flux
  // range plus the one extra face in -i.
  const auto rp = la::AddHalo<recon_halo>(r);
  auto qp = la::GetPerPointScratch<Real>(rp);
  auto qm = la::GetPerPointScratch<Real>(rp);
  auto pv = la::make_pack_view(rp, pack);

  // Producer: reconstruct over the halo-extended range.
  la::inner(rp, [&](auto f) {
    qp(f) = reconstruct_plus (pv(q(), f - dx1), pv(q(), f), pv(q(), f + dx1));
    qm(f) = reconstruct_minus(pv(q(), f - dx1), pv(q(), f), pv(q(), f + dx1));
  });
  r.TeamBarrier();  // producer must finish before the consumer reads (no-op where serial)

  // Consumer: one Riemann solve per face over the base range.
  auto fv = la::make_flux_pack_view(r, pack, X1DIR);
  la::inner(r, [&](auto f) {
    fv(q(), f) = riemann(qp(f - dx1), qm(f));
  });
});
\end{lstlisting}

Now the instantiations:
\begin{itemize}
\item \code{IndexSpace<loop_tag::bovi, inner_tag::logical_flat>} (or
  \code{loop_tag::bvoi}) recovers the CPU form. \code{outer} walks chunks of the
  block; \code{inner} sweeps a contiguous span; \code{GetPerPointScratch} resolves
  to team scratch (Kokkos) or a host bump-arena buffer (raw), sized to the
  halo-extended producer range; the \code{TeamBarrier} separates producer and
  consumer. The reconstructed states are stored once and reused. Choosing
  \code{bvoi} instead keeps the whole block live per slice --- more scratch,
  often faster (Section~\ref{sec:looptags}).
\item \code{IndexSpace<loop_tag::boiv, ...>} recovers the GPU form. \code{outer}
  is a flat grid over all $(\text{block},k,j,i)$, scratch is a tiny per-thread
  stack buffer, and the \code{TeamBarrier} is a no-op. The producer and consumer
  run back-to-back on the same thread, so the reconstructed states are effectively
  recomputed per face --- exactly Section~\ref{sec:gpu}. Note that here the word
  \code{inner} is something of a lie: for the base (non-halo) range the
  ``\code{inner} loop'' is a single logical cell --- a one-iteration loop, i.e.\
  not really a loop at all --- and even for the halo-extended range it runs only
  over the handful of halo offsets, not a sweep. Writing it as \code{inner}
  regardless is what lets the \emph{same} kernel text also compile to the
  hierarchical forms, where \code{inner} genuinely is a sweep; \code{boiv} is the
  degenerate limit in which that sweep collapses to one point.
\end{itemize}

One line in the body deserves emphasis, because it is where a large part of the
CPU/GPU performance difference is actually realized:
\code{make_pack_view(rp, pack)} adapts to the chosen structure. For the
hierarchical, flat-indexed paths (the CPU form) it pulls raw pointers out of the
\code{SparsePack} at construction, so \code{pv(q(), f)} compiles to an
\code{a[idx]}-style access the vectorizer can see through --- recovering the
raw-pointer access of Section~\ref{sec:cpu} without the kernel spelling any
pointer arithmetic. For the pointwise / coordinate paths
(\code{boiv}, \code{logical_coords}) it is instead a thin wrapper that forwards
each access straight to the sparse pack, which is what the GPU form wants. The
kernel text is identical either way; the view quietly does the right thing for the
tag. Section~\ref{sec:views} returns to this.

The flux kernel is the demanding case that motivated the design --- it exercises
scratch, halos, and a producer/consumer split all at once --- but it is not the
only intended use. The same \code{outer}/\code{inner} form is meant to be the way
ordinary block-structured kernels are written too, including ones that touch no
scratch and involve no stencil at all: a pointwise equation-of-state call, a
per-cell source term, a copy or a rescale. Such a kernel is simply the flux
example with the producer loop and halo removed --- one \code{inner} over the base
range --- and it still benefits from compile-time structure selection and the
pack-view access described above. The goal is a single idiom for block kernels,
with the reconstruction/flux machinery as the elaborate end of a spectrum whose
simple end is an unadorned loop.

The rest of this document explains each piece used above.

% =====================================================================
\section{Index spaces and coordinates}
\label{sec:indexspace}
% =====================================================================

\emph{[Draft scaffold --- to be expanded.]}

An \code{IndexSpace<loop_tag, inner_tag, backend>} carries two indexers for a
block: a \emph{logical} indexer describing the interior $(k,j,i)$ cells the kernel
must cover, and a \emph{memory} indexer describing the full allocated extent
including ghost cells. The distinction matters because halos and \code{memory}
traversals touch cells outside the logical set, and because \emph{degenerate}
directions (a symmetry direction, extent~1, no ghosts) must be distinguished from
thin active directions. Cover here: construction (the extent-based test
constructor vs.\ the \code{md}-based constructor), \code{GetLogicalIndexer} /
\code{GetMemoryIndexer}, flat-vs-coordinate conversion, and \code{GetNdim} (keyed
off the memory extents, consistent with \code{GetDelta}).

% =====================================================================
\section{Loop tags: where the variable level sits}
\label{sec:looptags}
% =====================================================================

\emph{[Draft scaffold --- to be expanded.]}

The \code{loop_tag} selects where the user's ``variable'' ($v$) loop --- any
per-block work between the block level and the innermost $kji$ traversal, most
often a loop over fields/components --- sits in the
$\text{block}\to\text{outer}\to\text{inner}$ hierarchy, and this fixes the
generated loop structure. Emphasize the chunking/fusion picture: \code{bovi} and
\code{boiv} \emph{chunk} (and can fuse) the $kji$ space in \code{outer}, while
\code{bvoi} keeps the whole block in one slice.
\begin{description}
\item[\code{bvoi}] $v$ above the traversal; a $v$ iteration spans the whole $kji$
  space of a block. Mixed logical/memory path; halos/scratch are block-relative.
\item[\code{bovi}] $v$ inside the chunk \code{outer} loop, around \code{inner};
  each $v$ iteration covers one $kji$ chunk. The main contiguous-span path
  (the CPU form of \S\ref{sec:cpu}).
\item[\code{boiv}] $v$ innermost, at a single logical cell. The pointwise hot path
  (the GPU form of \S\ref{sec:gpu}); the \code{ninner}${=}1$ limit of \code{bovi}
  with its own code path.
\end{description}

% =====================================================================
\section{Inner tags: how a chunk is traversed}
\label{sec:innertags}
% =====================================================================

\emph{[Draft scaffold --- to be expanded.]}

\begin{description}
\item[\code{logical_flat}] Each logical cell once, body receives a flat integer for
  \code{var[idx]} access; requires a shared layout across fields; most vectorizable.
\item[\code{logical_coords}] Same coverage, body receives an \code{Index3};
  use when fields have different layouts (e.g.\ cell- vs.\ face-centered).
\item[\code{memory}] Contiguous memory span for the chunk, may include ghost
  cells; logical cells still touched exactly once. \code{boiv}${+}$\code{memory}
  is rejected at compile time.
\end{description}

% =====================================================================
\section{Backends}
\label{sec:backends}
% =====================================================================

\emph{[Draft scaffold --- to be expanded.]}

The third template parameter is \code{loop_backend::raw} (a host nest with
\code{#pragma omp simd}) or \code{loop_backend::kokkos} (Kokkos parallel
policies). It defaults to \code{default_loop_backend_v} --- \code{raw} when the
device execution space is the host space, \code{kokkos} otherwise. \code{outer} /
\code{inner} dispatch with \code{if constexpr}, so selection is zero-cost.

% =====================================================================
\section{Scratch}
\label{sec:scratch}
% =====================================================================

\emph{[Draft scaffold --- to be expanded.]}

Per-point scratch is registered on the \code{IndexSpace} at setup
(\code{AddPerPointScratch<T>()}, \code{<T,2,3>()} for a shaped block,
\code{<T,Halo>()} to size for a halo-extended range) and requested inside the
outer body (\code{GetPerPointScratch<T>(range)}). The returned object is
\emph{specialized per loop path and backend} behind a uniform indexing API:
\begin{itemize}
\item \code{boiv}: a tiny fixed-size \emph{stack} buffer (a few cells).
\item raw backend: a non-owning slice of a per-thread \emph{bump arena} --- a
  region allocator that hands out memory by advancing a pointer and reclaims
  everything at once when the outer iteration resets
  (\href{https://en.wikipedia.org/wiki/Region-based_memory_management}{region-based
  allocation}).
\item Kokkos backend: Kokkos team scratch (shared memory).
\end{itemize}
As with raw hierarchical parallelism, call \code{range.TeamBarrier()} between a
producer inner loop and a consumer that reads values other threads wrote.

% =====================================================================
\section{Halos}
\label{sec:halos}
% =====================================================================

A \emph{halo} answers one question: \emph{if a consumer inner loop runs over a
logical point set $\S$, which neighboring produced values must also exist?} It is
a compile-time set of logical offsets that extends the point set a producer loop
must cover, so that a later consumer over $\S$ can read produced values at its
neighbors. This is the abstraction of the ``two loops over nearly the same range''
pattern from Section~\ref{sec:cpu}.

Formally, for a halo with offsets $\{h_1, h_2, \dots\}$ (the identity
$\mathbf{0}$ is always implicitly included),
\[
  \texttt{AddHalo}\langle h_1, h_2, \dots\rangle(\S)
  \;=\; \S \,\cup\, \{p + h_1 : p \in \S\} \,\cup\, \{p + h_2 : p \in \S\}
  \,\cup\, \dots
\]
In the flux kernel, the consumer reads \code{qp(f - dx1)}: a produced value one
cell in $-i$. So the producer must cover $\S \cup \{p - \hat{\imath}\}$, named by
\code{halo::minus_i_t}.

Figure~\ref{fig:halo} shows the geometry for one- and two-offset halos and for
several relationships between the halo direction and the inner-chunk layout.

\begin{figure}[t]
\centering
\begin{tikzpicture}[scale=0.62, every node/.style={font=\small}]
  % ============================================================
  % Visual convention:
  % blue    = previous inner range S_{m-1}
  % orange  = representative inner range S_m
  % green   = next inner range S_{m+1}
  % purple  = later inner range S_{m+2}
  % red x   = points visited by S_m.AddHalo<...>()
  % red dashed outline = shifted halo cells h(S_m)
  %
  % The drawn grid is the halo-extended logical domain.
  % ============================================================

  % ============================================================
  % Panel A: +j halo, ninner > ni
  % Original logical ni=6, nj=6. Extended grid has nj+1 rows.
  % ============================================================
  \begin{scope}[yshift=0cm]
    \def\ni{6}
    \def\njext{7}

    \node[font=\bfseries, anchor=west] at (0,7.8)
      {(a) $+j$ halo with $n_{\mathrm{inner}}=10>n_i=6$};

    % Halo-extended logical grid
    \draw[step=1, gray!45, very thin] (0,0) grid (\ni,\njext);

    % Original-domain inner ranges
    \fill[blue!18] (0,0) rectangle (6,1);
    \fill[blue!18] (0,1) rectangle (4,2);

    \fill[orange!40] (4,1) rectangle (6,2);
    \fill[orange!40] (0,2) rectangle (6,3);
    \fill[orange!40] (0,3) rectangle (2,4);

    \fill[green!25] (2,3) rectangle (6,4);
    \fill[green!25] (0,4) rectangle (6,5);

    % Mark globally added halo row
    \fill[red!8] (0,6) rectangle (6,7);
    \node[red!70!black, font=\scriptsize] at (3,6.5) {global $+j$ halo row};

    \draw[step=1, gray!60, very thin] (0,0) grid (\ni,\njext);

    % Boundary for S_m
    \draw[orange!85!black, very thick]
      (4,1) -- (6,1) -- (6,3) -- (2,3) -- (2,4) -- (0,4) -- (0,2) -- (4,2) -- cycle;

    % Shifted +j halo cells h(S_m)
    \draw[red!80!black, very thick, dashed]
      (4,2) -- (6,2) -- (6,4) -- (2,4) -- (2,5) -- (0,5) -- (0,3) -- (4,3) -- cycle;

    % S_m union shift_j(S_m)
    \foreach \i/\j in {
      4/1,5/1,
      0/2,1/2,2/2,3/2,4/2,5/2,
      0/3,1/3,2/3,3/3,4/3,5/3,
      0/4,1/4
    } {
      \node[red!80!black, font=\bfseries\large] at (\i+0.5,\j+0.5) {$\times$};
    }

    \node[blue!60!black] at (3,0.5) {$S_{m-1}$};
    \node[orange!85!black] at (3,2.5) {$S_m$};
    \node[green!45!black] at (3,4.5) {$S_{m+1}$};

    \node[anchor=west] at (6.25,0) {$i$};
    \node[anchor=south] at (0,7.1) {$j$};

    \node[align=left, anchor=west] at (7.4,5.2)
      {$S_m.\mathrm{AddHalo}\langle +j\rangle$\\
       $= S_m \cup \{p+\hat{\jmath}:p\in S_m\}$};

    % Per-point dependency inset
    \begin{scope}[xshift=8.3cm, yshift=1.4cm]
      \node[font=\small\bfseries, anchor=west] at (-0.2,2.45)
        {per-point dependency};

      \fill[orange!45] (0,0) rectangle (1,1);
      \draw[orange!85!black, very thick] (0,0) rectangle (1,1);
      \draw[red!80!black, very thick, dashed] (0,1) rectangle (1,2);

      \draw[red!80!black, thick] (0.5,0.5) -- (0.5,1.5);
      \fill[orange!85!black] (0.5,0.5) circle (0.08);
      \node[red!80!black, font=\bfseries\large] at (0.5,1.5) {$\times$};

      \node[anchor=east] at (-0.15,0.5) {$p$};
      \node[anchor=east, red!80!black] at (-0.15,1.5) {$p+\hat{\jmath}$};
    \end{scope}
  \end{scope}

  % ============================================================
  % Panel B: +i halo, ninner > ni
  % Original logical ni=6, nj=6. Extended grid has ni+1 columns.
  % ============================================================
  \begin{scope}[yshift=-9.5cm]
    \def\niext{7}
    \def\nj{6}

    \node[font=\bfseries, anchor=west] at (0,6.8)
      {(b) $+i$ halo with $n_{\mathrm{inner}}=10>n_i=6$};

    % Halo-extended logical grid
    \draw[step=1, gray!45, very thin] (0,0) grid (\niext,\nj);

    % Original-domain inner ranges
    \fill[blue!18] (0,0) rectangle (6,1);
    \fill[blue!18] (0,1) rectangle (4,2);

    \fill[orange!40] (4,1) rectangle (6,2);
    \fill[orange!40] (0,2) rectangle (6,3);
    \fill[orange!40] (0,3) rectangle (2,4);

    \fill[green!25] (2,3) rectangle (6,4);
    \fill[green!25] (0,4) rectangle (6,5);

    % Mark globally added halo column
    \fill[red!8] (6,0) rectangle (7,6);
    \node[red!70!black, rotate=90, font=\scriptsize] at (6.5,3)
      {global $+i$ halo column};

    \draw[step=1, gray!60, very thin] (0,0) grid (\niext,\nj);

    % Boundary for S_m
    \draw[orange!85!black, very thick]
      (4,1) -- (6,1) -- (6,3) -- (2,3) -- (2,4) -- (0,4) -- (0,2) -- (4,2) -- cycle;

    % Shifted +i halo cells h(S_m), geometrically shifted right.
    \draw[red!80!black, very thick, dashed]
      (5,1) -- (7,1) -- (7,3) -- (2,3) -- (2,4) -- (1,4) -- (1,2) -- (5,2) -- cycle;

    % S_m union shift_i(S_m)
    \foreach \i/\j in {
      4/1,5/1,6/1,
      0/2,1/2,2/2,3/2,4/2,5/2,6/2,
      0/3,1/3,2/3
    } {
      \node[red!80!black, font=\bfseries\large] at (\i+0.5,\j+0.5) {$\times$};
    }

    \node[blue!60!black] at (3,0.5) {$S_{m-1}$};
    \node[orange!85!black] at (3,2.5) {$S_m$};
    \node[green!45!black] at (4,3.5) {$S_{m+1}$};

    \node[anchor=west] at (7.25,0) {$i$};
    \node[anchor=south] at (0,6.1) {$j$};

    \node[align=left, anchor=west] at (8.2,4.75)
      {$S_m.\mathrm{AddHalo}\langle +i\rangle$\\
       $= S_m \cup \{p+\hat{\imath}:p\in S_m\}$};

    % Per-point dependency inset
    \begin{scope}[xshift=9.1cm, yshift=1.45cm]
      \node[font=\small\bfseries, anchor=west] at (-0.2,1.45)
        {per-point dependency};

      \fill[orange!45] (0,0) rectangle (1,1);
      \draw[orange!85!black, very thick] (0,0) rectangle (1,1);
      \draw[red!80!black, very thick, dashed] (1,0) rectangle (2,1);

      \draw[red!80!black, thick] (0.5,0.5) -- (1.5,0.5);
      \fill[orange!85!black] (0.5,0.5) circle (0.08);
      \node[red!80!black, font=\bfseries\large] at (1.5,0.5) {$\times$};

      \node[anchor=north] at (0.5,-0.15) {$p$};
      \node[anchor=north, red!80!black] at (1.5,-0.15) {$p+\hat{\imath}$};
    \end{scope}
  \end{scope}

  % ============================================================
  % Panel C: +j halo, ninner < ni
  % Original logical ni=8, nj=5. Extended grid has nj+1 rows.
  % ============================================================
  \begin{scope}[yshift=-18.1cm]
    \def\ni{8}
    \def\njext{6}

    \node[font=\bfseries, anchor=west] at (0,6.8)
      {(c) $+j$ halo with $n_{\mathrm{inner}}=4<n_i=8$};

    % Halo-extended logical grid
    \draw[step=1, gray!45, very thin] (0,0) grid (\ni,\njext);

    % Each range has length 4.
    \fill[blue!18]    (0,1) rectangle (4,2); % S_{m-1}
    \fill[orange!40]  (4,1) rectangle (8,2); % S_m
    \fill[green!25]   (0,2) rectangle (4,3); % S_{m+1}
    \fill[purple!20]  (4,2) rectangle (8,3); % S_{m+2}

    % Mark globally added halo row
    \fill[red!8] (0,5) rectangle (8,6);
    \node[red!70!black, font=\scriptsize] at (4,5.5) {global $+j$ halo row};

    \draw[step=1, gray!60, very thin] (0,0) grid (\ni,\njext);

    % Boundary for S_m
    \draw[orange!85!black, very thick]
      (4,1) rectangle (8,2);

    % Shifted +j halo cells h(S_m), coinciding with S_{m+2}
    \draw[red!80!black, very thick, dashed]
      (4,2) rectangle (8,3);

    % S_m union shift_j(S_m)
    \foreach \i/\j in {
      4/1,5/1,6/1,7/1,
      4/2,5/2,6/2,7/2
    } {
      \node[red!80!black, font=\bfseries\large] at (\i+0.5,\j+0.5) {$\times$};
    }

    \node[blue!60!black] at (2,1.5) {$S_{m-1}$};
    \node[orange!85!black] at (6,1.5) {$S_m$};
    \node[green!45!black] at (2,2.5) {$S_{m+1}$};
    \node[purple!60!black] at (6,2.5) {$S_{m+2}$};

    \node[anchor=west] at (8.25,0) {$i$};
    \node[anchor=south] at (0,6.1) {$j$};

    \node[align=left, anchor=west] at (9.5,4.6)
      {$S_m.\mathrm{AddHalo}\langle +j\rangle$\\
       $= S_m \cup \{p+\hat{\jmath}:p\in S_m\}$\\[2pt]
       {The shifted copy lands in $S_{m+2}$.}};

    % Per-point dependency inset, moved lower to avoid annotation overlap
    \begin{scope}[xshift=10.4cm, yshift=0.65cm]
      \node[font=\small\bfseries, anchor=west] at (-0.2,2.45)
        {per-point dependency};

      \fill[orange!45] (0,0) rectangle (1,1);
      \draw[orange!85!black, very thick] (0,0) rectangle (1,1);
      \draw[red!80!black, very thick, dashed] (0,1) rectangle (1,2);

      \draw[red!80!black, thick] (0.5,0.5) -- (0.5,1.5);
      \fill[orange!85!black] (0.5,0.5) circle (0.08);
      \node[red!80!black, font=\bfseries\large] at (0.5,1.5) {$\times$};

      \node[anchor=east] at (-0.15,0.5) {$p$};
      \node[anchor=east, red!80!black] at (-0.15,1.5) {$p+\hat{\jmath}$};
    \end{scope}
  \end{scope}
\end{tikzpicture}
\caption{A halo extends the logical point set an inner range visits. Orange is a
representative inner range $S_m$; the red dashed outline is its shifted copy under
the halo offset; red crosses mark the union
$\texttt{AddHalo}(S_m) = S_m \cup \text{shift}(S_m, h)$. Panels (a,b) show $+j$ and
$+i$ halos when a chunk spans more than one row; panel (c) shows a $+j$ halo when
chunks are single sub-rows, so the shifted copy lands in a later chunk. Insets show
the single-point dependency that the halo encodes.}
\label{fig:halo}
\end{figure}

% ---------------------------------------------------------------
\subsection{A halo is not a reconstruction stencil}
\label{sec:halo-not-stencil}
% ---------------------------------------------------------------

This distinction is the single most common point of confusion, so it is worth
stating sharply. The \emph{reconstruction stencil} is the set of \emph{input}
cells read to compute \emph{one} reconstructed value; in the flux kernel it is
wide (\code{q(f-2*dx1)} through \code{q(f+dx1)} for \code{reconstruct_plus}). The
\emph{halo} is the set of \emph{extra output faces} the producer must fill so the
consumer's neighbor reads are satisfied; here it is a single cell, $-\hat{\imath}$.

The two are unrelated in size. The stencil width is a property of the
reconstruction operator and lives entirely inside the producer body (the
arguments to \code{reconstruct_plus}). The halo is a property of the
producer/consumer \emph{coupling} --- how far the consumer reaches into produced
scratch --- and is expressed in the range the producer iterates. A wide stencil
with a one-cell halo is the common case: the reconstruction reads five cells, but
the flux loop only ever reads the produced state one face away.

% ---------------------------------------------------------------
\subsection{Offsets, \texttt{GetDelta}, and sign conventions}
\label{sec:getdelta}
% ---------------------------------------------------------------

\emph{[Draft scaffold --- to be expanded.]}

\code{GetDelta(dir)} returns an offset object \code{dx1} representing a
\emph{$+1$} logical step in direction \code{dir}, usable in the body's index
arithmetic: \code{f + dx1}, \code{f - dx1}, \code{f + 2*dx1}. It is signed through
the coefficient; \code{f - dx1} is the step in $-\hat{\imath}$. Cover here: why the
offset is an object rather than a raw integer (it carries the memory stride so it
composes with flat, \code{Index3}, and memory-offset index forms), and how it
relates to the halo offsets --- crucially, \code{GetDelta} zeroes the component of
a degenerate direction, which is the hook that connects to
Section~\ref{sec:reduced}.

% ---------------------------------------------------------------
\subsection{Halos in reduced dimensionality}
\label{sec:reduced}
% ---------------------------------------------------------------

A halo offset that points into a direction the problem does not resolve names a
shifted copy of the inner range that does not exist. A $k$-directed halo in a 2D
($x_1$--$x_2$) problem is the canonical case: there is no $k\pm1$ plane to shift
into. The abstraction must not let such an offset extend the index space, size
scratch, or drive a loop past the real cells.

\paragraph{DROP.} In a run where some direction is degenerate, the abstraction
keeps only the offsets whose components in every degenerate direction are zero,
and discards the rest. Because the halo offsets are stored sorted
lexicographically by $(dk, dj, di)$ and the degenerate directions are always the
most-significant sort keys ($k$, then $j$, in the Parthenon convention where
degenerate directions are the higher ones), the kept offsets form a single
contiguous run $[\text{begin}, \text{end})$. The identity is always kept. This
contiguity is what makes the reduction cheap: the span-merge and pointwise
enumeration that build a halo range simply iterate the sub-range instead of the
full offset list, and the sortedness they rely on is preserved.

\paragraph{PROJECT, and why DROP is usually enough.} The semantically correct
operation is not DROP but \emph{PROJECT}: replace each offset by its projection
onto the active directions (zero the degenerate components). Under the assumption
that a removed dimension is a genuine \emph{translational symmetry} direction ---
every field is constant along it --- PROJECT is exact. If the consumer reads a
produced value at $p+h$, symmetry makes that value equal to the one at the
projected point $\text{proj}(p+h)$, so producing the projected copy delivers
precisely what the consumer needs, and the reduced-dimension run agrees with the
fully-resolved run.

DROP and PROJECT agree exactly when the halo is \emph{closed under projection} ---
when every offset's projection is itself already a declared offset. This holds for
every physical stencil we expect: dense corner boxes (e.g.\ the $2\times2$ face-flux
corners of constrained-transport EMFs), single-direction extensions, and standard
3/5/7-point stencils. It fails only for sparse, L-shaped, or bare-diagonal halos ---
an offset like $(-1,0,2)$ present without its projection $(0,0,2)$ --- which do not
arise in the schemes this abstraction targets.

\paragraph{The compile-time guard.} Because DROP silently under-produces when the
halo is \emph{not} projection-closed, \code{AddHalo} enforces closure with a
\code{static_assert}. The check is dimension-independent: it tests closure under
both degenerations Parthenon can produce (2D, with $k$ inactive; and 1D, with $k$
and $j$ inactive), so a halo that passes is safe in any run regardless of runtime
dimensionality. A halo that fails is rejected at compile time with a message
directing the author to the fix: add the missing projection point. Doing so is
free of side effects --- it produces one extra scratch value that, by projection
closure, no consumer distinguishes from a value it already needed --- and it makes
DROP coincide with PROJECT for that halo. Implementing PROJECT proper (which
requires de-duplicating the offsets that collapse onto the same projected point) is
a deliberately deferred extension; until a genuinely non-closed halo is needed in
practice, the guard plus the one-line workaround covers every case.

% =====================================================================
\section{Pack, variable, and flux views}
\label{sec:views}
% =====================================================================

\emph{[Draft scaffold --- to be expanded.]}

Views adapt a \code{SparsePack} to the current loop contract so variable access
uses the same index forms as the body (flat \code{int}, \code{Index3}, or explicit
$k,j,i$):
\begin{itemize}
\item \code{make_pack_view} (non-sparse vars), \code{make_sparse_pack_view}
  (sparse vars at a sparse index), \code{make_var_view} (single variable).
\item \code{make_flux_pack_view(range, pack, dir)} / \code{make_flux_view} --- flux
  counterparts. Note the deliberate difference from ordinary Parthenon packs: a
  flux view contains \emph{only} fluxes and \emph{only} for the one direction
  requested at construction.
\end{itemize}
For \code{logical_coords} these are thin wrappers over the pack; for the other
inner tags they pull out raw pointers at construction, which can significantly aid
vectorization.

% =====================================================================
\section{Contracts and invariants}
\label{sec:contracts}
% =====================================================================

\emph{[Draft scaffold --- to be expanded. Mirror the invariant summary from
\code{LOOP_ABSTRACTION_CONTRACTS.md}.]}

The invariants every path must satisfy:
\begin{enumerate}
\item Every loop pattern touches every logical cell exactly once.
\item \code{logical_flat} and \code{logical_coords} must not touch halo/ghost cells.
\item \code{memory} may touch ghost cells but must still touch every logical cell
  exactly once.
\item raw and Kokkos backends agree on the contract for the same loop pattern and
  body signature.
\end{enumerate}

\end{document}
